\section{Interactive constructs and graceful degradation}

pandoc-carve runs \textbf{no} renderer extensions - it maps the parsed
core AST to Pandoc. So every script-dependent construct (tabs,
code-group, spoiler, mermaid, math) degrades to its static form on the
way to LaTeX, Typst, DOCX, and the rest. Content and structure survive;
only the interaction is dropped. This file is the worked proof - compare
it against the generated \texttt{interactive.tex} / \texttt{.typ} /
\texttt{.native} beside it.

\subsection{Tabs}

A \texttt{{[}label{]}} on each panel is a grouping identifier. With a
tabs extension it becomes a clickable tab button; here, with no
extension, it degrades to a visible caption so a print reader can still
tell the panels apart.

\textbf{Installation}

Run \texttt{npm\ install\ @markup-carve/pandoc-carve}.

\textbf{Usage}

Call the \emph{converter} on your \texttt{.crv} source.

\subsection{Code group}

Same grouping mechanism, code panels instead of prose. Each fence keeps
its language (so syntax highlighting survives every writer); the panels
stack.

\begin{Shaded}
\begin{Highlighting}[]
\ImportTok{export} \ImportTok{default}\NormalTok{ \{ }\DataTypeTok{port}\OperatorTok{:} \DecValTok{3000}\NormalTok{ \}}\OperatorTok{;}
\end{Highlighting}
\end{Shaded}

\begin{Shaded}
\begin{Highlighting}[]
\FunctionTok{\{} \DataTypeTok{"port"}\FunctionTok{:} \DecValTok{3000} \FunctionTok{\}}
\end{Highlighting}
\end{Shaded}

\subsection{Spoiler}

A quoted title survives as a caption; the ``hide until revealed''
interaction is meaningless offline, so the body is simply shown.

\textbf{Answer}

The hidden \emph{answer} is 42.

\subsection{Disclosure}

\textbf{Show details}

A disclosure carries its title and body through as a labelled block.

\subsection{Mermaid and math}

The diagram source is never lost - it stays a fenced code block that a
Markdown host can re-render, and a build step can pre-render to an image
for PDF.

\begin{Shaded}
\begin{Highlighting}[]
\NormalTok{graph TD; A[Carve] {-}{-}\textgreater{} B[Pandoc]; B {-}{-}\textgreater{} C[LaTeX];}
\end{Highlighting}
\end{Shaded}

Inline math keeps its source too:
\(\sum_{i=1}^{n} i = \frac{n(n+1)}{2}\).
